\textbf{Phần 1. Trắc nghiệm lựa chọn}
\setcounterpageref{ex}{0}

\begin{ex} %[5D12V73]
Hai ô tô xuất phát cùng một lúc từ hai tỉnh A và B cách nhau $240\text{ km}$, đi ngược chiều nhau. Sau $2$ giờ $30$ phút thì hai xe gặp nhau. Biết vận tốc của ô tô đi từ A lớn hơn vận tốc của ô tô đi từ B là $8\text{ km/h}$. Vận tốc của ô tô đi từ A là
\choice
{$44\text{ km/h}$}
{\True $52\text{ km/h}$}
{$50\text{ km/h}$}
{$56\text{ km/h}$}
\loigiai{
Đổi: $2\text{ giờ } 30\text{ phút} = 2{,}5\text{ giờ}$.\\
Tổng vận tốc của hai xe là:
$$240 : 2{,}5 = 96\text{ (km/h)}$$
Vận tốc của ô tô đi từ A là:
$$(96 + 8) : 2 = 52\text{ (km/h)}$$
\textbf{Đáp án của câu hỏi là B}
}
\end{ex}

\begin{ex} %[5D12V73]
Một người đi xe đạp từ A với vận tốc $12\text{ km/h}$. Sau $1$ giờ $30$ phút, một người đi xe máy cũng xuất phát từ A đuổi theo xe đạp với vận tốc $36\text{ km/h}$. Thời gian kể từ lúc xe máy xuất phát đến khi đuổi kịp xe đạp là
\choice
{$30\text{ phút}$}
{\True $45\text{ phút}$}
{$1\text{ giờ}$}
{$1\text{ giờ } 15\text{ phút}$}
\loigiai{
Đổi: $1\text{ giờ } 30\text{ phút} = 1{,}5\text{ giờ}$.\\
Quãng đường người đi xe đạp đi trước là:
$$12 \times 1{,}5 = 18\text{ (km)}$$
Hiệu vận tốc giữa hai xe là:
$$36 - 12 = 24\text{ (km/h)}$$
Thời gian xe máy đuổi kịp xe đạp là:
$$18 : 24 = 0{,}75\text{ (giờ)} = 45\text{ (phút)}$$
\textbf{Đáp án của câu hỏi là B}
}
\end{ex}

\begin{ex} %[5D12V73]
Một chiếc ca nô xuôi dòng từ A đến B mất $2$ giờ $30$ phút và ngược dòng từ B về A mất $3$ giờ $45$ phút. Biết vận tốc của dòng nước là $3\text{ km/h}$. Vận tốc thực của ca nô khi nước lặng là
\choice
{$12\text{ km/h}$}
{\True $15\text{ km/h}$}
{$18\text{ km/h}$}
{$21\text{ km/h}$}
\loigiai{
Đổi: $2\text{ giờ } 30\text{ phút} = 2{,}5\text{ giờ}$; $3\text{ giờ } 45\text{ phút} = 3{,}75\text{ giờ}$.\\
Tỉ số thời gian xuôi dòng và ngược dòng là:
$$2{,}5 : 3{,}75 = \dfrac{2}{3}$$
Vì quãng đường không đổi nên tỉ số vận tốc xuôi dòng và ngược dòng là $\dfrac{3}{2}$.\\
Hiệu vận tốc xuôi dòng và ngược dòng bằng $2$ lần vận tốc dòng nước:
$$3 \times 2 = 6\text{ (km/h)}$$
Vận tốc xuôi dòng của ca nô là:
$$6 : (3 - 2) \times 3 = 18\text{ (km/h)}$$
Vận tốc thực của ca nô khi nước lặng là:
$$18 - 3 = 15\text{ (km/h)}$$
\textbf{Đáp án của câu hỏi là B}
}
\end{ex}

\begin{ex} %[5D12V73]
Khoảng cách giữa hai bến sông A và B ở bài toán trên là bao nhiêu ki-lô-mét?
\choice
{$40\text{ km}$}
{\True $45\text{ km}$}
{$50\text{ km}$}
{$60\text{ km}$}
\loigiai{
Vận tốc xuôi dòng của ca nô là $18\text{ km/h}$ và thời gian xuôi dòng là $2{,}5\text{ giờ}$.\\
Khoảng cách giữa hai bến A và B là:
$$18 \times 2{,}5 = 45\text{ (km)}$$
\textbf{Đáp án của câu hỏi là B}
}
\end{ex}

\begin{ex} %[5D12V73]
Một người đi xe máy từ A đến B với vận tốc $40\text{ km/h}$, sau đó từ B quay trở về A với vận tốc $60\text{ km/h}$. Vận tốc trung bình của người đó trên toàn bộ hành trình cả đi lẫn về là
\choice
{$50\text{ km/h}$}
{\True $48\text{ km/h}$}
{$45\text{ km/h}$}
{$52\text{ km/h}$}
\loigiai{
Giả sử quãng đường AB dài $120\text{ km}$.\\
Thời gian lúc đi là:
$$120 : 40 = 3\text{ (giờ)}$$
Thời gian lúc về là:
$$120 : 60 = 2\text{ (giờ)}$$
Tổng quãng đường cả đi lẫn về là:
$$120 \times 2 = 240\text{ (km)}$$
Tổng thời gian cả đi lẫn về là:
$$3 + 2 = 5\text{ (giờ)}$$
Vận tốc trung bình trên cả hành trình là:
$$240 : 5 = 48\text{ (km/h)}$$
\textbf{Đáp án của câu hỏi là B}
}
\end{ex}

\begin{ex} %[5D12V73]
Một đoàn tàu hỏa chạy ngang qua một cột điện bên đường mất $12\text{ giây}$ và chạy qua một cây cầu dài $360\text{ m}$ mất $30\text{ giây}$. Vận tốc của đoàn tàu là
\choice
{$54\text{ km/h}$}
{\True $72\text{ km/h}$}
{$60\text{ km/h}$}
{$80\text{ km/h}$}
\loigiai{
Thời gian đoàn tàu đi được chiều dài của cây cầu ($360\text{ m}$) là:
$$30 - 12 = 18\text{ (giây)}$$
Vận tốc của đoàn tàu là:
$$360 : 18 = 20\text{ (m/s)}$$
Đổi: $20\text{ m/s} = 20 \times 3{,}6 = 72\text{ (km/h)}$.\\
\textbf{Đáp án của câu hỏi là B}
}
\end{ex}

\begin{ex} %[5D12V73]
Chiều dài của đoàn tàu hỏa ở bài toán trên là bao nhiêu mét?
\choice
{$200\text{ m}$}
{\True $240\text{ m}$}
{$280\text{ m}$}
{$300\text{ m}$}
\loigiai{
Đoàn tàu chạy ngang qua cột điện mất $12\text{ giây}$ với vận tốc $20\text{ m/s}$.\\
Chiều dài của đoàn tàu là:
$$20 \times 12 = 240\text{ (m)}$$
\textbf{Đáp án của câu hỏi là B}
}
\end{ex}

\begin{ex} %[5D12V73]
Quãng đường AB dài $195\text{ km}$. Lúc $7$ giờ, một xe máy đi từ A đến B với vận tốc $30\text{ km/h}$. Đến $7$ giờ $30$ phút, một ô tô đi từ B về A với vận tốc $60\text{ km/h}$. Hai xe gặp nhau vào thời điểm nào?
\choice
{$9$ giờ}
{\True $9$ giờ $30$ phút}
{$9$ giờ $15$ phút}
{$10$ giờ}
\loigiai{
Từ $7$ giờ đến $7$ giờ $30$ phút là $30\text{ phút} = 0{,}5\text{ giờ}$.\\
Quãng đường xe máy đi trước là:
$$30 \times 0{,}5 = 15\text{ (km)}$$
Quãng đường còn lại khi ô tô xuất phát là:
$$195 - 15 = 180\text{ (km)}$$
Tổng vận tốc của hai xe là:
$$30 + 60 = 90\text{ (km/h)}$$
Thời gian để hai xe gặp nhau kể từ lúc ô tô xuất phát là:
$$180 : 90 = 2\text{ (giờ)}$$
Hai xe gặp nhau lúc:
$$7\text{ giờ } 30\text{ phút} + 2\text{ giờ} = 9\text{ giờ } 30\text{ phút}$$
\textbf{Đáp án của câu hỏi là B}
}
\end{ex}

\begin{ex} %[5D12V73]
Một ô tô khởi hành từ A lúc $8$ giờ để đi đến B dài $150\text{ km}$ với vận tốc $50\text{ km/h}$. Trên đường đi, ô tô dừng lại nghỉ $20$ phút. Ô tô đến B lúc mấy giờ?
\choice
{$11$ giờ}
{\True $11$ giờ $20$ phút}
{$11$ giờ $30$ phút}
{$10$ giờ $50$ phút}
\loigiai{
Thời gian ô tô chạy trên đường là:
$$150 : 50 = 3\text{ (giờ)}$$
Tổng thời gian di chuyển cả nghỉ là:
$$3\text{ giờ} + 20\text{ phút} = 3\text{ giờ } 20\text{ phút}$$
Ô tô đến B lúc:
$$8\text{ giờ} + 3\text{ giờ } 20\text{ phút} = 11\text{ giờ } 20\text{ phút}$$
\textbf{Đáp án của câu hỏi là B}
}
\end{ex}

\begin{ex} %[5D12V73]
Một người đi bộ từ nhà lên huyện với vận tốc $4\text{ km/h}$, khi từ huyện về nhà người đó đi bằng xe đạp với vận tốc $12\text{ km/h}$. Thời gian lúc đi nhiều hơn thời gian lúc về là $2$ giờ. Quãng đường từ nhà lên huyện dài bao nhiêu ki-lô-mét?
\choice
{$10\text{ km}$}
{\True $12\text{ km}$}
{$16\text{ km}$}
{$18\text{ km}$}
\loigiai{
Tỉ số vận tốc giữa lúc đi và lúc về là:
$$4 : 12 = \dfrac{1}{3}$$
Vì quãng đường không đổi nên thời gian tỉ lệ nghịch với vận tốc, thời gian lúc đi gấp $3$ lần thời gian lúc về.\\
Hiệu số phần bằng nhau là:
$$3 - 1 = 2\text{ (phần)}$$
Thời gian lúc về là:
$$2 : 2 \times 1 = 1\text{ (giờ)}$$
Quãng đường từ nhà lên huyện là:
$$12 \times 1 = 12\text{ (km)}$$
\textbf{Đáp án của câu hỏi là B}
}
\end{ex}

\begin{ex} %[5D12V73]
Một đường đua vòng quanh bờ hồ dài $400\text{ m}$. Hai vận động viên An và Bình cùng xuất phát tại một vạch và chạy cùng chiều nhau. Vận tốc của An là $5\text{ m/s}$, của Bình là $3\text{ m/s}$. Sau bao lâu kể từ lúc xuất phát thì An đuổi kịp Bình lần đầu tiên?
\choice
{$2$ phút}
{\True $3\text{ phút } 20\text{ giây}$}
{$4$ phút}
{$2\text{ phút } 40\text{ giây}$}
\loigiai{
Khi An đuổi kịp Bình lần đầu tiên thì An đã chạy nhiều hơn Bình đúng $1$ vòng hồ ($400\text{ m}$).\\
Hiệu vận tốc của hai người là:
$$5 - 3 = 2\text{ (m/s)}$$
Thời gian để An đuổi kịp Bình là:
$$400 : 2 = 200\text{ (giây)} = 3\text{ phút } 20\text{ giây}$$
\textbf{Đáp án của câu hỏi là B}
}
\end{ex}

\begin{ex} %[5D12V73]
Lúc $6$ giờ sáng, một xe tải khởi hành từ A đi về B với vận tốc $40\text{ km/h}$. Đến $7$ giờ $15$ phút, một xe ô tô con cũng xuất phát từ A đuổi theo xe tải với vận tốc $60\text{ km/h}$. Điểm ô tô con đuổi kịp xe tải cách A quãng đường dài bao nhiêu ki-lô-mét?
\choice
{$120\text{ km}$}
{\True $150\text{ km}$}
{$160\text{ km}$}
{$180\text{ km}$}
\loigiai{
Thời gian xe tải đi trước là:
$$7\text{ giờ } 15\text{ phút} - 6\text{ giờ} = 1\text{ giờ } 15\text{ phút} = 1{,}25\text{ giờ}$$
Quãng đường xe tải đi trước là:
$$40 \times 1{,}25 = 50\text{ (km)}$$
Hiệu vận tốc giữa ô tô con và xe tải là:
$$60 - 40 = 20\text{ (km/h)}$$
Thời gian ô tô con đuổi kịp xe tải là:
$$50 : 20 = 2{,}5\text{ (giờ)}$$
Điểm gặp nhau cách A là:
$$60 \times 2{,}5 = 150\text{ (km)}$$
\textbf{Đáp án của câu hỏi là B}
}
\end{ex}

\textbf{Phần 2. Trắc nghiệm đúng sai}
\setcounterpageref{ex}{0}

\begin{ex} %[5D12V73]
Hai địa điểm A và B cách nhau $160\text{ km}$. Lúc $7$ giờ $30$ phút, một người đi xe máy khởi hành từ A về B với vận tốc $35\text{ km/h}$ và một ô tô khởi hành từ B về A với vận tốc $45\text{ km/h}$. Xét tính đúng sai của các khẳng định sau:
\choiceTF
{\True Tổng vận tốc của xe máy và ô tô là $80\text{ km/h}$.}
{\True Sau $2$ giờ kể từ lúc cùng khởi hành thì hai xe gặp nhau.}
{Hai xe gặp nhau vào thời điểm $9$ giờ $15$ phút sáng cùng ngày.}
{\True Vị trí hai xe gặp nhau cách địa điểm A một khoảng cách là $70\text{ km}$.}
\loigiai{
\begin{itemchoice}
    \itemch Đúng. Tổng vận tốc: $35 + 45 = 80\text{ km/h}$.
    \itemch Đúng. Thời gian gặp nhau: $160 : 80 = 2\text{ giờ}$.
    \itemch Sai. Thời điểm gặp nhau: $7\text{ giờ } 30\text{ phút} + 2\text{ giờ} = 9\text{ giờ } 30\text{ phút} \neq 9\text{ giờ } 15\text{ phút}$.
    \itemch Đúng. Điểm gặp cách A: $35 \times 2 = 70\text{ km}$.
\end{itemchoice}
}
\end{ex}

\begin{ex} %[5D12V73]
Một ca nô đi xuôi dòng từ bến A đến bến B hết $2$ giờ và đi ngược dòng từ B về A hết $3$ giờ. Biết khoảng cách giữa hai bến A và B là $60\text{ km}$. Xét tính đúng sai của các khẳng định sau:
\choiceTF
{\True Vận tốc của ca nô khi xuôi dòng là $30\text{ km/h}$.}
{\True Vận tốc của ca nô khi ngược dòng là $20\text{ km/h}$.}
{Vận tốc của dòng nước là $10\text{ km/h}$.}
{\True Vận tốc thực của ca nô khi nước lặng là $25\text{ km/h}$.}
\loigiai{
\begin{itemchoice}
    \itemch Đúng. Vận tốc xuôi dòng: $60 : 2 = 30\text{ km/h}$.
    \itemch Đúng. Vận tốc ngược dòng: $60 : 3 = 20\text{ km/h}$.
    \itemch Sai. Vận tốc dòng nước: $(30 - 20) : 2 = 5\text{ km/h} \neq 10\text{ km/h}$.
    \itemch Đúng. Vận tốc thực của ca nô: $30 - 5 = 25\text{ km/h}$.
\end{itemchoice}
}
\end{ex}

\textbf{Phần 3. Trả lời ngắn}
\setcounterpageref{ex}{0}

\begin{ex} %[5D12V73]
Hai ô tô khởi hành cùng một lúc từ hai địa điểm cách nhau $225\text{ km}$, đi ngược chiều nhau với vận tốc lần lượt là $42\text{ km/h}$ và $48\text{ km/h}$. Hỏi sau bao nhiêu phút kể từ lúc xuất phát thì hai ô tô gặp nhau?
\shortans{150}
\loigiai{
Tổng vận tốc của hai ô tô là:
$$42 + 48 = 90\text{ (km/h)}$$
Thời gian hai xe gặp nhau là:
$$225 : 90 = 2{,}5\text{ (giờ)} = 150\text{ (phút)}$$
\textbf{Đáp số:} $150$
}
\end{ex}

\begin{ex} %[5D12V73]
Một người đi xe đạp với vận tốc $15\text{ km/h}$ trong thời gian $1$ giờ $24$ phút. Tính quãng đường người đó đã đi được theo ki-lô-mét.
\shortans{21}
\loigiai{
Đổi: $1\text{ giờ } 24\text{ phút} = 1{,}4\text{ giờ}$.\\
Quãng đường người đó đi được là:
$$15 \times 1{,}4 = 21\text{ (km)}$$
\textbf{Đáp số:} $21$
}
\end{ex}

\begin{ex} %[5D12V73]
Một người đi xe máy từ A với vận tốc $36\text{ km/h}$. Sau $40$ phút, một ô tô cũng xuất phát từ A đuổi theo với vận tốc $60\text{ km/h}$. Hỏi sau bao nhiêu phút kể từ lúc ô tô xuất phát thì ô tô đuổi kịp xe máy?
\shortans{60}
\loigiai{
Đổi: $40\text{ phút} = \dfrac{2}{3}\text{ giờ}$.\\
Quãng đường xe máy đi trước là:
$$36 \times \dfrac{2}{3} = 24\text{ (km)}$$
Hiệu vận tốc giữa hai xe là:
$$60 - 36 = 24\text{ (km/h)}$$
Thời gian ô tô đuổi kịp xe máy là:
$$24 : 24 = 1\text{ (giờ)} = 60\text{ (phút)}$$
\textbf{Đáp số:} $60$
}
\end{ex}

\begin{ex} %[5D12V73]
Một đoàn tàu hỏa dài $150\text{ m}$ chạy qua một cây cầu dài $450\text{ m}$ với vận tốc $54\text{ km/h}$. Hỏi thời gian để đoàn tàu vượt qua hoàn toàn cây cầu đó là bao nhiêu giây?
\shortans{40}
\loigiai{
Đổi: $54\text{ km/h} = 54 : 3{,}6 = 15\text{ m/s}$.\\
Quãng đường đoàn tàu cần đi để vượt qua hết cây cầu là:
$$150 + 450 = 600\text{ (m)}$$
Thời gian đoàn tàu đi qua hết cây cầu là:
$$600 : 15 = 40\text{ (giây)}$$
\textbf{Đáp số:} $40$
}
\end{ex}

\begin{ex} %[5D12V73]
Một ca nô đi xuôi dòng khúc sông từ A đến B với vận tốc $28\text{ km/h}$ và đi ngược dòng từ B về A với vận tốc $20\text{ km/h}$. Tính vận tốc của dòng nước theo $\text{km/h}$.
\shortans{4}
\loigiai{
Vận tốc của dòng nước là:
$$(28 - 20) : 2 = 4\text{ (km/h)}$$
\textbf{Đáp số:} $4$
}
\end{ex}

\begin{ex} %[5D12V73]
Bác An đi bộ thể dục quanh một công viên hình chữ nhật có chiều dài $160\text{ m}$, chiều rộng $90\text{ m}$. Bác đi được đúng $4$ vòng với vận tốc đều $4\text{ km/h}$. Hỏi bác An đã đi trong bao nhiêu phút?
\shortans{30}
\loigiai{
Chu vi công viên là:
$$(160 + 90) \times 2 = 500\text{ (m)}$$
Quãng đường bác An đi được sau $4$ vòng là:
$$500 \times 4 = 2\,000\text{ (m)} = 2\text{ (km)}$$
Thời gian bác An đi là:
$$2 : 4 = 0{,}5\text{ (giờ)} = 30\text{ (phút)}$$
\textbf{Đáp số:} $30$
}
\end{ex}
